% Build after collecting the figures into build/for-paper/
% (python3 build/sim_paper_figures/collect_for_paper.py from the repository root):
%   (cd docs/slides && latexmk -pdf -interaction=nonstopmode -halt-on-error \
%      -outdir=../../build/for-paper 2026_10_13_adc_performance.tex)
% Shared FRIDA Beamer style. Start each slide deck with % Shared FRIDA Beamer style. Start each slide deck with % Shared FRIDA Beamer style. Start each slide deck with \input{style.tex}.
\documentclass[aspectratio=169,compress]{beamer}
\usepackage[T1]{fontenc}
\usepackage{lmodern,graphicx,xcolor,amsmath,booktabs,array,tikz,listings,etoolbox}

% Decks are built from docs/slides; figures live beside docs and in docs/images.
\graphicspath{{../}{./}}

\definecolor{NordBody}{HTML}{4C566A}
\definecolor{NordBlue}{HTML}{5E81AC}
\definecolor{NordRed}{HTML}{BF616A}
\definecolor{NordGreen}{HTML}{A3BE8C}
\definecolor{NordOrange}{HTML}{D08770}
\definecolor{NordPurple}{HTML}{B48EAD}
\definecolor{NordPanel}{HTML}{ECEFF4}

\newcommand{\framesource}{}
\BeforeBeginEnvironment{frame}{\gdef\framesource{}}
\newcommand{\slidesource}[1]{\gdef\framesource{#1}}
\setbeamertemplate{footline}{%
  \hbox{%
    \begin{beamercolorbox}[wd=.12\paperwidth,ht=2.7ex,dp=1.2ex,leftskip=1em]{page number in head/foot}%
      \usebeamerfont{page number in head/foot}\insertframenumber\,/\,\inserttotalframenumber
    \end{beamercolorbox}%
    \begin{beamercolorbox}[wd=.88\paperwidth,ht=2.7ex,dp=1.2ex,right,rightskip=1em]{normal text}%
      \fontsize{5.5}{6.5}\selectfont\ifx\framesource\empty\else Sources:\enspace\framesource\fi
    \end{beamercolorbox}%
  }%
}
\beamertemplatenavigationsymbolsempty
\setbeameroption{hide notes}

% Opt in for decks with sections; shorter decks keep the full frame height.
\newcommand{\slideprogressnavigation}{%
  \useoutertheme[subsection=false]{miniframes}
  \setbeamercolor{section in head/foot}{fg=NordBlue,bg=white}
  \setbeamercolor{mini frame}{fg=NordBlue,bg=white}
  \setbeamercolor{lower separation line head}{bg=NordPanel}
  \setbeamerfont{section in head/foot}{size=\tiny}

  \newif\ifnavbeforecurrent
  \pretocmd\insertnavigation{\navbeforecurrenttrue}{}{}

  \newcommand{\fridanavigationdot}{%
    \begin{pgfpicture}{0pt}{0pt}{0.1cm}{0.1cm}
      \pgfpathcircle{\pgfpoint{0.05cm}{0.05cm}}{0.05cm}
      \ifnavbeforecurrent
        \pgfusepath{fill,stroke}
      \else
        \pgfusepath{stroke}
      \fi
    \end{pgfpicture}%
  }

  \defbeamertemplate*{mini frame}{frida progress}{%
    \begin{pgfpicture}{0pt}{0pt}{0.1cm}{0.1cm}
      \pgfpathcircle{\pgfpoint{0.05cm}{0.05cm}}{0.05cm}
      \pgfusepath{fill,stroke}
    \end{pgfpicture}%
    \global\navbeforecurrentfalse
  }
  \defbeamertemplate*{mini frame in current section}{frida progress}{\fridanavigationdot}
  \defbeamertemplate*{mini frame in current subsection}{frida progress}{\fridanavigationdot}
  \defbeamertemplate*{mini frame in other section}{frida progress}{\fridanavigationdot}
  \defbeamertemplate*{mini frame in other subsection}{frida progress}{\fridanavigationdot}
  \setbeamertemplate{mini frame}[frida progress]
  \setbeamertemplate{mini frame in current section}[frida progress]
  \setbeamertemplate{mini frame in current subsection}[frida progress]
  \setbeamertemplate{mini frame in other section}[frida progress]
  \setbeamertemplate{mini frame in other subsection}[frida progress]
}

\setbeamerfont{frametitle}{size=\large}
\setbeamerfont{framesubtitle}{size=\small}
\setbeamercolor{normal text}{fg=NordBody}
\setbeamercolor{title}{fg=NordBlue}
\setbeamercolor{frametitle}{fg=NordBlue}
\setbeamercolor{framesubtitle}{fg=NordBody}
\setbeamercolor{structure}{fg=NordBlue}
\setbeamercolor{itemize/enumerate body}{fg=NordBody}
\setbeamercolor{itemize/enumerate subbody}{fg=NordBody}
\setbeamercolor{itemize/enumerate subsubbody}{fg=NordBody}
\setbeamercolor{itemize item}{fg=NordBlue}
\setbeamercolor{itemize subitem}{fg=NordBlue}
\setbeamercolor{itemize subsubitem}{fg=NordBlue}
\setbeamercolor{enumerate item}{fg=NordBlue}
\setbeamercolor{enumerate subitem}{fg=NordBlue}
\setbeamercolor{enumerate subsubitem}{fg=NordBlue}
\setbeamercolor{block title}{fg=NordBlue,bg=NordPanel}
\setbeamercolor{block body}{fg=NordBody,bg=NordPanel!55}
\hypersetup{colorlinks=true,allcolors=NordBlue}

\lstset{%
  basicstyle=\ttfamily\tiny,
  keywordstyle=\color{NordBlue}\bfseries,
  commentstyle=\color{NordBody}\itshape,
  stringstyle=\color{NordOrange},
  showstringspaces=false,
  breaklines=true,
  frame=single,
  backgroundcolor=\color{NordPanel!55},
  rulecolor=\color{NordBody!50},
  xleftmargin=2pt,
  xrightmargin=2pt,
  aboveskip=4pt,
  belowskip=4pt,
}

\newcommand{\placeholder}[2]{%
  \begin{tikzpicture}
    \node[minimum width=#1,minimum height=#2,draw=NordBlue!60,rounded corners=2pt,fill=NordPanel,align=center,text=NordBody]
      {Placeholder\\[-1pt]\scriptsize image not yet available};
  \end{tikzpicture}%
}
\newcommand{\maybeimage}[2][]{%
  \IfFileExists{#2}{\includegraphics[#1]{#2}}{%
    \IfFileExists{../#2}{\includegraphics[#1]{#2}}{\placeholder{0.95\linewidth}{0.30\textheight}}%
  }%
}

\author{Kennedy Caisley}
\institute{University of Bonn}
.
\documentclass[aspectratio=169,compress]{beamer}
\usepackage[T1]{fontenc}
\usepackage{lmodern,graphicx,xcolor,amsmath,booktabs,array,tikz,listings,etoolbox}

% Decks are built from docs/slides; figures live beside docs and in docs/images.
\graphicspath{{../}{./}}

\definecolor{NordBody}{HTML}{4C566A}
\definecolor{NordBlue}{HTML}{5E81AC}
\definecolor{NordRed}{HTML}{BF616A}
\definecolor{NordGreen}{HTML}{A3BE8C}
\definecolor{NordOrange}{HTML}{D08770}
\definecolor{NordPurple}{HTML}{B48EAD}
\definecolor{NordPanel}{HTML}{ECEFF4}

\newcommand{\framesource}{}
\BeforeBeginEnvironment{frame}{\gdef\framesource{}}
\newcommand{\slidesource}[1]{\gdef\framesource{#1}}
\setbeamertemplate{footline}{%
  \hbox{%
    \begin{beamercolorbox}[wd=.12\paperwidth,ht=2.7ex,dp=1.2ex,leftskip=1em]{page number in head/foot}%
      \usebeamerfont{page number in head/foot}\insertframenumber\,/\,\inserttotalframenumber
    \end{beamercolorbox}%
    \begin{beamercolorbox}[wd=.88\paperwidth,ht=2.7ex,dp=1.2ex,right,rightskip=1em]{normal text}%
      \fontsize{5.5}{6.5}\selectfont\ifx\framesource\empty\else Sources:\enspace\framesource\fi
    \end{beamercolorbox}%
  }%
}
\beamertemplatenavigationsymbolsempty
\setbeameroption{hide notes}

% Opt in for decks with sections; shorter decks keep the full frame height.
\newcommand{\slideprogressnavigation}{%
  \useoutertheme[subsection=false]{miniframes}
  \setbeamercolor{section in head/foot}{fg=NordBlue,bg=white}
  \setbeamercolor{mini frame}{fg=NordBlue,bg=white}
  \setbeamercolor{lower separation line head}{bg=NordPanel}
  \setbeamerfont{section in head/foot}{size=\tiny}

  \newif\ifnavbeforecurrent
  \pretocmd\insertnavigation{\navbeforecurrenttrue}{}{}

  \newcommand{\fridanavigationdot}{%
    \begin{pgfpicture}{0pt}{0pt}{0.1cm}{0.1cm}
      \pgfpathcircle{\pgfpoint{0.05cm}{0.05cm}}{0.05cm}
      \ifnavbeforecurrent
        \pgfusepath{fill,stroke}
      \else
        \pgfusepath{stroke}
      \fi
    \end{pgfpicture}%
  }

  \defbeamertemplate*{mini frame}{frida progress}{%
    \begin{pgfpicture}{0pt}{0pt}{0.1cm}{0.1cm}
      \pgfpathcircle{\pgfpoint{0.05cm}{0.05cm}}{0.05cm}
      \pgfusepath{fill,stroke}
    \end{pgfpicture}%
    \global\navbeforecurrentfalse
  }
  \defbeamertemplate*{mini frame in current section}{frida progress}{\fridanavigationdot}
  \defbeamertemplate*{mini frame in current subsection}{frida progress}{\fridanavigationdot}
  \defbeamertemplate*{mini frame in other section}{frida progress}{\fridanavigationdot}
  \defbeamertemplate*{mini frame in other subsection}{frida progress}{\fridanavigationdot}
  \setbeamertemplate{mini frame}[frida progress]
  \setbeamertemplate{mini frame in current section}[frida progress]
  \setbeamertemplate{mini frame in current subsection}[frida progress]
  \setbeamertemplate{mini frame in other section}[frida progress]
  \setbeamertemplate{mini frame in other subsection}[frida progress]
}

\setbeamerfont{frametitle}{size=\large}
\setbeamerfont{framesubtitle}{size=\small}
\setbeamercolor{normal text}{fg=NordBody}
\setbeamercolor{title}{fg=NordBlue}
\setbeamercolor{frametitle}{fg=NordBlue}
\setbeamercolor{framesubtitle}{fg=NordBody}
\setbeamercolor{structure}{fg=NordBlue}
\setbeamercolor{itemize/enumerate body}{fg=NordBody}
\setbeamercolor{itemize/enumerate subbody}{fg=NordBody}
\setbeamercolor{itemize/enumerate subsubbody}{fg=NordBody}
\setbeamercolor{itemize item}{fg=NordBlue}
\setbeamercolor{itemize subitem}{fg=NordBlue}
\setbeamercolor{itemize subsubitem}{fg=NordBlue}
\setbeamercolor{enumerate item}{fg=NordBlue}
\setbeamercolor{enumerate subitem}{fg=NordBlue}
\setbeamercolor{enumerate subsubitem}{fg=NordBlue}
\setbeamercolor{block title}{fg=NordBlue,bg=NordPanel}
\setbeamercolor{block body}{fg=NordBody,bg=NordPanel!55}
\hypersetup{colorlinks=true,allcolors=NordBlue}

\lstset{%
  basicstyle=\ttfamily\tiny,
  keywordstyle=\color{NordBlue}\bfseries,
  commentstyle=\color{NordBody}\itshape,
  stringstyle=\color{NordOrange},
  showstringspaces=false,
  breaklines=true,
  frame=single,
  backgroundcolor=\color{NordPanel!55},
  rulecolor=\color{NordBody!50},
  xleftmargin=2pt,
  xrightmargin=2pt,
  aboveskip=4pt,
  belowskip=4pt,
}

\newcommand{\placeholder}[2]{%
  \begin{tikzpicture}
    \node[minimum width=#1,minimum height=#2,draw=NordBlue!60,rounded corners=2pt,fill=NordPanel,align=center,text=NordBody]
      {Placeholder\\[-1pt]\scriptsize image not yet available};
  \end{tikzpicture}%
}
\newcommand{\maybeimage}[2][]{%
  \IfFileExists{#2}{\includegraphics[#1]{#2}}{%
    \IfFileExists{../#2}{\includegraphics[#1]{#2}}{\placeholder{0.95\linewidth}{0.30\textheight}}%
  }%
}

\author{Kennedy Caisley}
\institute{University of Bonn}
.
\documentclass[aspectratio=169,compress]{beamer}
\usepackage[T1]{fontenc}
\usepackage{lmodern,graphicx,xcolor,amsmath,booktabs,array,tikz,listings,etoolbox}

% Decks are built from docs/slides; figures live beside docs and in docs/images.
\graphicspath{{../}{./}}

\definecolor{NordBody}{HTML}{4C566A}
\definecolor{NordBlue}{HTML}{5E81AC}
\definecolor{NordRed}{HTML}{BF616A}
\definecolor{NordGreen}{HTML}{A3BE8C}
\definecolor{NordOrange}{HTML}{D08770}
\definecolor{NordPurple}{HTML}{B48EAD}
\definecolor{NordPanel}{HTML}{ECEFF4}

\newcommand{\framesource}{}
\BeforeBeginEnvironment{frame}{\gdef\framesource{}}
\newcommand{\slidesource}[1]{\gdef\framesource{#1}}
\setbeamertemplate{footline}{%
  \hbox{%
    \begin{beamercolorbox}[wd=.12\paperwidth,ht=2.7ex,dp=1.2ex,leftskip=1em]{page number in head/foot}%
      \usebeamerfont{page number in head/foot}\insertframenumber\,/\,\inserttotalframenumber
    \end{beamercolorbox}%
    \begin{beamercolorbox}[wd=.88\paperwidth,ht=2.7ex,dp=1.2ex,right,rightskip=1em]{normal text}%
      \fontsize{5.5}{6.5}\selectfont\ifx\framesource\empty\else Sources:\enspace\framesource\fi
    \end{beamercolorbox}%
  }%
}
\beamertemplatenavigationsymbolsempty
\setbeameroption{hide notes}

% Opt in for decks with sections; shorter decks keep the full frame height.
\newcommand{\slideprogressnavigation}{%
  \useoutertheme[subsection=false]{miniframes}
  \setbeamercolor{section in head/foot}{fg=NordBlue,bg=white}
  \setbeamercolor{mini frame}{fg=NordBlue,bg=white}
  \setbeamercolor{lower separation line head}{bg=NordPanel}
  \setbeamerfont{section in head/foot}{size=\tiny}

  \newif\ifnavbeforecurrent
  \pretocmd\insertnavigation{\navbeforecurrenttrue}{}{}

  \newcommand{\fridanavigationdot}{%
    \begin{pgfpicture}{0pt}{0pt}{0.1cm}{0.1cm}
      \pgfpathcircle{\pgfpoint{0.05cm}{0.05cm}}{0.05cm}
      \ifnavbeforecurrent
        \pgfusepath{fill,stroke}
      \else
        \pgfusepath{stroke}
      \fi
    \end{pgfpicture}%
  }

  \defbeamertemplate*{mini frame}{frida progress}{%
    \begin{pgfpicture}{0pt}{0pt}{0.1cm}{0.1cm}
      \pgfpathcircle{\pgfpoint{0.05cm}{0.05cm}}{0.05cm}
      \pgfusepath{fill,stroke}
    \end{pgfpicture}%
    \global\navbeforecurrentfalse
  }
  \defbeamertemplate*{mini frame in current section}{frida progress}{\fridanavigationdot}
  \defbeamertemplate*{mini frame in current subsection}{frida progress}{\fridanavigationdot}
  \defbeamertemplate*{mini frame in other section}{frida progress}{\fridanavigationdot}
  \defbeamertemplate*{mini frame in other subsection}{frida progress}{\fridanavigationdot}
  \setbeamertemplate{mini frame}[frida progress]
  \setbeamertemplate{mini frame in current section}[frida progress]
  \setbeamertemplate{mini frame in current subsection}[frida progress]
  \setbeamertemplate{mini frame in other section}[frida progress]
  \setbeamertemplate{mini frame in other subsection}[frida progress]
}

\setbeamerfont{frametitle}{size=\large}
\setbeamerfont{framesubtitle}{size=\small}
\setbeamercolor{normal text}{fg=NordBody}
\setbeamercolor{title}{fg=NordBlue}
\setbeamercolor{frametitle}{fg=NordBlue}
\setbeamercolor{framesubtitle}{fg=NordBody}
\setbeamercolor{structure}{fg=NordBlue}
\setbeamercolor{itemize/enumerate body}{fg=NordBody}
\setbeamercolor{itemize/enumerate subbody}{fg=NordBody}
\setbeamercolor{itemize/enumerate subsubbody}{fg=NordBody}
\setbeamercolor{itemize item}{fg=NordBlue}
\setbeamercolor{itemize subitem}{fg=NordBlue}
\setbeamercolor{itemize subsubitem}{fg=NordBlue}
\setbeamercolor{enumerate item}{fg=NordBlue}
\setbeamercolor{enumerate subitem}{fg=NordBlue}
\setbeamercolor{enumerate subsubitem}{fg=NordBlue}
\setbeamercolor{block title}{fg=NordBlue,bg=NordPanel}
\setbeamercolor{block body}{fg=NordBody,bg=NordPanel!55}
\hypersetup{colorlinks=true,allcolors=NordBlue}

\lstset{%
  basicstyle=\ttfamily\tiny,
  keywordstyle=\color{NordBlue}\bfseries,
  commentstyle=\color{NordBody}\itshape,
  stringstyle=\color{NordOrange},
  showstringspaces=false,
  breaklines=true,
  frame=single,
  backgroundcolor=\color{NordPanel!55},
  rulecolor=\color{NordBody!50},
  xleftmargin=2pt,
  xrightmargin=2pt,
  aboveskip=4pt,
  belowskip=4pt,
}

\newcommand{\placeholder}[2]{%
  \begin{tikzpicture}
    \node[minimum width=#1,minimum height=#2,draw=NordBlue!60,rounded corners=2pt,fill=NordPanel,align=center,text=NordBody]
      {Placeholder\\[-1pt]\scriptsize image not yet available};
  \end{tikzpicture}%
}
\newcommand{\maybeimage}[2][]{%
  \IfFileExists{#2}{\includegraphics[#1]{#2}}{%
    \IfFileExists{../#2}{\includegraphics[#1]{#2}}{\placeholder{0.95\linewidth}{0.30\textheight}}%
  }%
}

\author{Kennedy Caisley}
\institute{University of Bonn}

\slideprogressnavigation

\newcommand{\figpath}{../../build/for-paper}
\newcommand{\plotfig}[2][0.80]{%
  \centering
  \includegraphics[width=0.98\linewidth,height=#1\textheight,keepaspectratio]{\figpath/#2.pdf}%
}

\title{ADC performance plots}
\subtitle{Behavioral simulations of the FRIDA-2 three-layer SAR ADC}
\date{13 October 2026}

\begin{document}

\begin{frame}
  \titlepage
\end{frame}

\begin{frame}{Simulation setup}
  \centering
  {\small
  \begin{tabular}{@{}ll@{}}
    \toprule
    Array & FRIDA-2 three-layer, radix 1.7: 16 capacitors, 17 decisions, 12-bit output \\
    Capacitances & PEX table, with Pelgrom mismatch $\sigma_C = A\sqrt{C\cdot\mathrm{fF}}$ \\
    Mismatch & Pelgrom model; coefficient $A = 0.85\,\%$ (Tripathi) or $0.10\,\%$ (Harpe) \\
    Noise & kT/C (1210\,fF), comparator 130\,\textmu V, 2\,ps aperture jitter (assumed) \\
    Settling & single pole, $\tau = 0.2$\,ns (assumed) \\
    Energy & logic, drivers and comparator fitted to the SPICE PEX transient \\
    Decodings & ideal design weights, PEX-extracted weights, calibrations 1--3 \\
    \bottomrule
  \end{tabular}
  }
  \slidesource{\texttt{flow/adc/behavioral.py}; \texttt{flow/scans/scan\_behavioral.py}}
\end{frame}

\section[Linearity]{Static linearity}

\begin{frame}{INL and DNL: Pelgrom mismatch, Tripathi coefficient}
  \plotfig[0.74]{01_inl_dnl_tripathi}
  \par\smallskip
  {\small $A = 0.85\,\%\sqrt{\mathrm{fF}}$. Ideal weights: peak $|$INL$|$ 5.5\,LSB, systematic. PEX-extracted weights: 1.4\,LSB, the random mismatch.}
  \slidesource{Behavioral held-out ramp, 10\,MS/s, draw 1; $A$ from Tripathi \& Murmann (2014)}
\end{frame}

\begin{frame}{INL and DNL: Pelgrom mismatch, Harpe coefficient}
  \plotfig[0.74]{01_inl_dnl_harpe}
  \par\smallskip
  {\small $A = 0.10\,\%\sqrt{\mathrm{fF}}$. The ideal-weight sawtooth is systematic, from the PEX capacitor ratios; PEX-extracted weights remove it.}
  \slidesource{Behavioral held-out ramp, 10\,MS/s, draw 1; $A$ estimated from Harpe's unit-length DAC (2018)}
\end{frame}

\begin{frame}{Code density: Pelgrom mismatch, two coefficients}
  \begin{columns}[T]
    \begin{column}{0.49\linewidth}
      \centering
      {\small $A = 0.85\,\%\sqrt{\mathrm{fF}}$ (Tripathi)}\par
      \includegraphics[width=\linewidth]{\figpath/01_code_density_tripathi.pdf}
    \end{column}
    \begin{column}{0.49\linewidth}
      \centering
      {\small $A = 0.10\,\%\sqrt{\mathrm{fF}}$ (Harpe)}\par
      \includegraphics[width=\linewidth]{\figpath/01_code_density_harpe.pdf}
    \end{column}
  \end{columns}
  \par\medskip
  \centering
  {\small Code count over the ideal count, averaged in 128 bins.}
  \slidesource{Behavioral held-out ramps, mismatch draw 1}
\end{frame}

\section[Speed]{Redundancy, speed and power}

\begin{frame}{Redundancy recovers a random first decision}
  \plotfig{02_trajectory_density}
  \slidesource{Behavioral, 2{,}000 conversions of a 0\,V input at 10\,MS/s, all noise sources}
\end{frame}

\begin{frame}{ENOB versus sample rate}
  \plotfig{03_enob_vs_rate}
  \par\smallskip
  {\small The high-rate roll-off follows the assumed 0.2\,ns CDAC settling.}
  \slidesource{Behavioral sine at $-0.5$\,dBFS, 8{,}192 conversions per point, PEX-extracted decoding}
\end{frame}

\begin{frame}{Supply power versus sample rate}
  \plotfig{04_power_vs_rate}
  \par\smallskip
  {\small About 27\,pJ per conversion; SAR logic dominates. Static power is about 2.3\,\textmu W.}
  \slidesource{Behavioral model, energies fitted to the FRIDA-2 three-layer SPICE PEX transient}
\end{frame}

\begin{frame}{Energy per conversion versus input}
  \plotfig{05_energy_vs_input}
  \slidesource{Behavioral, 129 DC inputs at 10\,MS/s}
\end{frame}

\section[Noise]{Noise and dynamic performance}

\begin{frame}{Input-referred noise spectral density}
  \plotfig{06_noise_spectrum}
  \par\smallskip
  {\small Total 175\,\textmu V RMS (0.47\,LSB): comparator 142\,\textmu V, sampling switch 62\,\textmu V.}
  \slidesource{Behavioral, 50\,mV DC input, 100{,}000 conversions at 10\,MS/s}
\end{frame}

\begin{frame}{ENOB versus input frequency}
  \plotfig{07_enob_vs_input_frequency}
  \slidesource{Behavioral sine at $-0.5$\,dBFS, nine input frequencies up to Nyquist at 2, 6 and 10\,MS/s}
\end{frame}

\begin{frame}{Output spectra at 10\,MS/s}
  \begin{columns}[T]
    \begin{column}{0.49\linewidth}
      \centering
      {\small Input near DC}\par
      \includegraphics[width=\linewidth]{\figpath/07_spectrum_10msps_low.pdf}
    \end{column}
    \begin{column}{0.49\linewidth}
      \centering
      {\small Input near Nyquist}\par
      \includegraphics[width=\linewidth]{\figpath/07_spectrum_10msps_nyquist.pdf}
    \end{column}
  \end{columns}
  \slidesource{Behavioral sine at $-0.5$\,dBFS, 8{,}192 conversions, PEX-extracted decoding}
\end{frame}

\section[Calibration]{Calibration}

\begin{frame}{Three calibration methods and their sources}
  \centering
  {\small
  \begin{tabular}{@{}>{\raggedright\arraybackslash}p{0.20\linewidth}>{\raggedright\arraybackslash}p{0.42\linewidth}>{\raggedright\arraybackslash}p{0.30\linewidth}@{}}
    \toprule
    Method & What it measures & Source \\
    \midrule
    Cal.\ 1: CDAC S-curves & Each capacitor's top-plate step, from comparator S-curves with one capacitor switched A$\to$B & FRIDA CDAC characterization scan (our own method) \\
    Cal.\ 2: least-squares fit & All weights at once, regressing BOUT on the ideal code of a known ramp (ridge, $w \geq 0$) & Generic foreground calibration with a known input \\
    Cal.\ 3: slow-ramp thresholds & Each decision's 50\,\% threshold after an all-0 and an all-1 prefix; their spacing is the weight & A.\ Hsu, MIT PhD dissertation 2013, Algorithm~I (Sec.~4.2) \\
    \bottomrule
  \end{tabular}
  }
  \par\medskip
  {\small All three are digital foreground calibrations: they learn BOUT weights once and leave the analog ADC unchanged.}
  \slidesource{\texttt{flow/analysis/calibration1.py}, \texttt{calibration2.py}, \texttt{calibration3.py}}
\end{frame}

\begin{frame}{Comparing the calibration methods}
  \begin{columns}[c]
    \begin{column}{0.55\linewidth}
      \centering
      \includegraphics[width=\linewidth,height=0.84\textheight,keepaspectratio]{\figpath/08_calibrations_inl_dnl.pdf}
    \end{column}
    \begin{column}{0.40\linewidth}
      {\footnotesize
      \begin{tabular}{@{}lrr@{}}
        \toprule
        Decoding & $|$INL$|$ & $|$DNL$|$ \\
        \midrule
        Ideal weights & 5.53 & 1.00 \\
        PEX-extracted & 1.37 & 0.96 \\
        Cal.\ 1, CDAC S-curves & 0.83 & 1.00 \\
        Cal.\ 2, least squares & 1.52 & 1.05 \\
        Cal.\ 3, Hsu thresholds & 1.06 & 0.79 \\
        \bottomrule
      \end{tabular}
      \par\medskip
      Peak errors in LSB on the held-out ramp. Calibrations 2 and 3 learn from a separate training ramp.
      }
    \end{column}
  \end{columns}
  \slidesource{Behavioral, Pelgrom $A = 0.85\,\%$ (Tripathi), draw 1}
\end{frame}

\begin{frame}{Learned weights}
  \plotfig{08_calibration_weights}
  \slidesource{Behavioral, Pelgrom $A = 0.85\,\%$ (Tripathi), draw 1}
\end{frame}

\end{document}
